%% Erläuterungen zu den Befehlen erfolgen unter
%% diesem Beispiel.

%% Article Template
\documentclass[]{scrartcl}


%% UTF8 Encoding
\usepackage[utf8]{inputenc}
\usepackage[T1]{fontenc}
%\usepackage[margin=5cm]{geometry}% http://ctan.org/pkg/geometry
\usepackage{lmodern}
\usepackage{subcaption}
\usepackage{setspace}
\usepackage[french]{babel}
\usepackage{csquotes}
\usepackage{pdfpages}
% appendix
\usepackage[titletoc]{appendix}

\setstretch{0.99}

%Tabellen mit fixen Breiten
\usepackage{tabularx}

%% Grafik
\usepackage{graphicx} 
\newcommand{\name}{Insta.edit}

%% Links im PDF
\usepackage{hyperref}
\usepackage{rotating}
\usepackage{lscape}
\usepackage{fancyhdr}
\usepackage{float}
\usepackage[titles]{tocloft}
\pagestyle{fancy}
\rhead{\includegraphics[scale=0.15]{img/insa-logo}}
\chead{\Large -SIE- }
\lhead{Hexanome 4222}

\title{Rapport Synthes\\
\subtitle{}
\author{PLD: SIE\\
Professeurs:\
\ \\
Version 1.0}
\date{16/10/2019}}


\begin{document}

\maketitle

\begin{figure}[h]
	\centering
  \includegraphics[width=0.5\textwidth]{img/insa-logo}
	\label{fig:logo}
\end{figure}


\begin{center}
  \begin{tabular}{ | l | r | }
    \hline
    \textbf{Corentin Laharotte}\\ \hline
    \textbf{Cédric Milinaire }\\ \hline
    \textbf{Felix Castillon}\\ \hline
    \textbf{Roxane Debord}\\ \hline
     \textbf{Grazia Giulia Ribbeni}\\ \hline
     \textbf{Ousmane Touat} \\ \hline
     \textbf{David Hamidovic} \\ \hline


  \end{tabular}
\end{center}

\thispagestyle{empty}
\pagebreak
\vspace*{10pt}
\tableofcontents
\listoffigures
\newpage

\section{Comprehension de l'entreprise}
\section{Comprehension des Processus}
Nous nous intéressons ici au processus de gestion des contrats de maintenance et services avec ses sous-processus qui démarre lorsqu'une opportunité d'avenant ou d'un contrat de service arrive.\\

Le premier sous processus "Offre et revue d'offre" concerne l'analyse de l'offre du point de vue des ressources existantes et nécessaires, des risques et de la faisabilité. Cela peut entraîner l'arrêt du processus de maintenance entier si l'étude de faisabilité n'est pas satisfaisant. Si au contraire l'étude de faisabilité est à bilan positif une étude complète de la solution à proposer est effectuée et validée (ou pas) en interne, puis elle est transmise au client.\\

Le processus "Réalisation des travaux induits" démarre lorsque des travaux supplémentaires sont détectés au cours d'une maintenance. Après l'analyse des travaux induits et des aspects contractuel deux options sont possibles. Si le contrat inclut la réalisation des travaux induits, ceux-ci sont chiffrés et réalisé en chiffrage contrôlé. Sinon un nouveau devis est nécessaire. Dans les deux cas, la facturation a lieu une fois que toutes les validations ont été faites et les travaux induits ont été conclus.


\section{Propositions de solutions}
\end{document}
